%% kapa: exam=1 round=22 year=2011 period=1 form=A subject=civil_law subject_ko=민법 seq=15 number=15
%% source: Q-Net 감정평가사 기출문제 (cst00309) artlSeq=5036790 file=extracted/1008/5036790/2011년 제22회 감정평가사 1차시험 기출문제/1교시.hwp
%% conversion: kordoc 4.24.51 -> tools/gen1.py (auto); review=proofread-sonnet; flags=none
\begin{question}{15}
전세권에 관한 설명으로 옳지 않은 것은? (다툼이 있으면 판례에 의함)
\begin{choices}
  \item 전세권자에게는 경매청구권이 있다.
  \item 전세금의 지급은 전세권 성립요소가 아니다.
  \item 전세권자는 전세권설정자에게 필요비상환을 청구할 수 없다.
  \item 전세권자가 목적물을 타인에게 임대한 경우, 임대하지 않았으면 면할 수 있었던 불가항력으로 인한 손해에 대하여도 책임을 부담한다.
  \item 전세권이 성립된 후 목적물의 소유권이 이전된 경우, 종전 소유자는 원칙적으로 전세권설정자의 지위를 상실하여 전세금반환의무를 면한다.
\end{choices}
\end{question}
\answer{2}
